% Luc Destombe --- licence CC BY-NC-SA 4.0
% https://creativecommons.org/licenses/by-nc-sa/4.0/deed.fr
% Fichier autonome : compiler avec pdflatex, deux fois (signets, nombre de pages).
\documentclass[a4paper,10pt]{article}
\makeatletter
\newif\iflycee@niveaux
\lycee@niveauxtrue
\newif\iflycee@palatino
\lycee@palatinotrue

\RequirePackage[utf8]{inputenc}
\RequirePackage[T1]{fontenc}
\RequirePackage[french]{babel}
\RequirePackage{etoolbox}

\newcommand{\ShorthandsOff}{\@ifundefined{shorthandoff}{}{\shorthandoff{;:!?}}}
\newcommand{\ShorthandsOn}{\@ifundefined{shorthandon}{}{\shorthandon{;:!?}}}
\ShorthandsOff
\AtBeginDocument{\ShorthandsOn}
\AtBeginEnvironment{tikzpicture}{\ShorthandsOff}

\RequirePackage{geometry}
\RequirePackage{fancyhdr}
\RequirePackage{lastpage}
\RequirePackage{multicol}
\RequirePackage{array}
\RequirePackage{enumitem}
\RequirePackage{xcolor}
\RequirePackage{needspace}

\RequirePackage{amsmath,amssymb}
\RequirePackage{mathpazo}
\DeclareMathSizes{10.95}{10.95}{8}{5.5}
\begingroup\catcode`\_=\active \catcode`\^=\active
\gdef_#1{\sb{\scriptscriptstyle #1}}
\gdef^#1{\sp{\scriptscriptstyle\lycee@moinsepais #1}}
\endgroup
\newcommand\lycee@moins{\mathbin{%
  \setbox\z@\hbox{$\m@th\scriptscriptstyle\lycee@moinsorig$}%
  \dimen@=.55\wd\z@ \advance\dimen@-.5pt
  \hbox to.75\wd\z@{\hss$\m@th\scriptscriptstyle\vcenter{\hbox{%
    \tikz\draw[line width=.5pt,line cap=round](0,0)--(\the\dimen@,0);}}$\hss}}}
\begingroup\lccode`\~=`\- \lowercase{\endgroup\let~}\lycee@moins
\newcommand\lycee@moinsepais{\mathcode`\-="8000 }
\AtBeginDocument{\mathchardef\lycee@moinsorig=\mathcode`\-\relax}
\AtBeginDocument{\catcode`\_=12 \mathcode`\_="8000
  \catcode`\^=12 \mathcode`\^="8000 }
\newcommand\lycee@exposants{%
  \fontdimen13\textfont2=.48\fontdimen6\textfont2
  \fontdimen14\textfont2=.48\fontdimen6\textfont2
  \fontdimen15\textfont2=.48\fontdimen6\textfont2
  \fontdimen13\scriptfont2=.48\fontdimen6\scriptfont2
  \fontdimen14\scriptfont2=.48\fontdimen6\scriptfont2
  \fontdimen15\scriptfont2=.48\fontdimen6\scriptfont2}
\every@math@size\expandafter{\the\every@math@size\lycee@exposants}
\newlength{\retraitcalculs}
\setlength{\retraitcalculs}{2em}

\RequirePackage{tikz}
\usetikzlibrary{arrows.meta,calc,babel,shapes.misc}
\RequirePackage[most]{tcolorbox}
\tcbuselibrary{skins,breakable}

\definecolor{coulDef}{HTML}{1F4E79}
\definecolor{coulProp}{HTML}{7030A0}
\definecolor{coulRem}{HTML}{C55A11}
\definecolor{coulEx}{HTML}{2E7D32}
\definecolor{coulExo}{HTML}{B03A2E}
\definecolor{coulPart}{HTML}{1F4E79}
\definecolor{coulMeth}{HTML}{1B6E4F}
\definecolor{coulCourbe}{HTML}{0B5ED7}
\definecolor{coulCourbeB}{HTML}{C00000}
\definecolor{coulCourbeC}{HTML}{2E7D32}
\definecolor{coulGrille}{HTML}{8C8C8C}
\definecolor{coulRep}{HTML}{1B6E4F}
\definecolor{coulBar}{HTML}{2E7D32}

\newcommand{\R}{\mathbb{R}}
\newcommand{\Rs}{\mathbb{R}^{*}}
\renewcommand{\le}{\leqslant}
\renewcommand{\ge}{\geqslant}
\renewcommand{\approx}{\simeq}
\newcommand{\vect}[1]{\overrightarrow{#1}}
\newcommand{\norme}[1]{\left\lVert #1 \right\rVert}
\newcommand{\intervalle}[2]{\left[#1\,;#2\right]}
\RequirePackage{cancel}
\renewcommand{\CancelColor}{\color{coulExo}}
\newcommand{\fois}[1]{\times{\color{coulExo}#1}}
\newcommand{\barre}[1]{\cancel{#1}}

\tikzset{grille/.style={coulGrille,line width=0.4pt}}
\newcommand{\quadrillage}[4]{\draw[grille,step=1] (#1,#2) grid (#3,#4);}
\tikzset{
  croix/.style={cross out,draw,line width=0.6pt,inner sep=0pt,minimum size=4pt},
  croixdroite/.style={croix,rotate=45,line width=0.8pt},
}
\newcommand{\pt}[3]{\path (#1) node[croix]{} node[#2,font=\small,inner sep=2.5pt]{$#3$};}

\newcommand{\lycee@repere}[4]{%
  \draw[grille,step=1] (#1,#2) grid (#3,#4);
  \draw[-{Stealth[length=2mm]},thick] (#1,0)--({#3+0.4},0) node[below left=-1pt]{$x$};
  \draw[-{Stealth[length=2mm]},thick] (0,#2)--(0,{#4+0.4}) node[below left=-1pt]{$y$};
  \node[below left,font=\scriptsize,inner sep=1.5pt] at (0,0) {$0$};}
\newcommand{\lycee@gradx}[1]{\foreach \k/\l in {#1}{%
  \draw (\k,0.07)--(\k,-0.07) node[below,font=\scriptsize,inner sep=1.5pt]{$\l$};}}
\newcommand{\lycee@grady}[1]{\foreach \k/\l in {#1}{%
  \draw (0.07,\k)--(-0.07,\k) node[left,font=\scriptsize,inner sep=1.5pt]{$\l$};}}
\newcommand{\lycee@intff}[2]{\left[#1\,;#2\right]}
\newcommand{\lycee@intfo}[2]{\left[#1\,;#2\right[}
\newcommand{\lycee@intof}[2]{\left]#1\,;#2\right]}
\newcommand{\lycee@intoo}[2]{\left]#1\,;#2\right[}
\newcommand{\lycee@ens}[1]{\left\{#1\right\}}
\AtBeginDocument{%
  \providecommand{\repere}{\lycee@repere}%
  \providecommand{\gradx}{\lycee@gradx}%
  \providecommand{\grady}{\lycee@grady}%
  \providecommand{\Cf}{\mathcal{C}\sb{\scriptscriptstyle f}}%
  \providecommand{\Cg}{\mathcal{C}\sb{\scriptscriptstyle g}}%
  \providecommand{\intff}{\lycee@intff}%
  \providecommand{\intfo}{\lycee@intfo}%
  \providecommand{\intof}{\lycee@intof}%
  \providecommand{\intoo}{\lycee@intoo}%
  \providecommand{\ens}{\lycee@ens}}

\newlist{questions}{enumerate}{3}
\setlist[questions,1]{label=\arabic*\textdegree),leftmargin=*,itemsep=2pt,topsep=3pt}
\setlist[questions,2]{label=\alph*),leftmargin=*,itemsep=1pt,topsep=2pt}
\setlist[questions,3]{label=\roman*),leftmargin=*,itemsep=1pt,topsep=1pt}
\setlist[itemize]{label=\textbullet,leftmargin=*,itemsep=2pt,topsep=3pt}

\newcommand{\besoinplace}[1]{\par\penalty\@M\begingroup
  \dimen@=#1\relax\dimen@ii=\pagegoal\advance\dimen@ii-\pagetotal
  \ifdim\dimen@>\dimen@ii\ifdim\dimen@ii>\z@\vfil\fi\break\fi\endgroup}

\newcommand{\lycee@pied}{\thepage/\pageref{LastPage}}
\newcommand{\entete}[2]{%
  \pagestyle{fancy}\fancyhf{}%
  \fancyhead[L]{\footnotesize\color{coulPart}\textbf{#1}}%
  \fancyhead[R]{\footnotesize\color{coulPart}\textbf{#2}}%
  \fancyfoot[C]{\footnotesize\lycee@pied}%
  \renewcommand{\headrulewidth}{0.4pt}%
  \renewcommand{\headrule}{\hbox to\headwidth{%
    \color{coulPart!40}\leaders\hrule height \headrulewidth\hfill}}}

\RequirePackage{ccicons}
\newcommand{\lycee@licence}{{\scriptsize\color{gray}%
  \copyright\ Luc Destombe --- \ccbyncsa\ CC BY-NC-SA 4.0}}
\newif\iflycee@licence \lycee@licencetrue
\newcommand{\sanslicence}{\lycee@licencefalse}
\AtBeginDocument{\iflycee@licence\AddToHookNext{shipout/foreground}{%
  \put(0,0){\rlap{\kern\dimexpr1in+\hoffset+\oddsidemargin+\textwidth\relax
    \llap{\raisebox{-\dimexpr1in+\voffset+\topmargin+\headheight+\headsep
      +\textheight+\footskip\relax}{\lycee@licence}}}}}\fi}

\newcommand{\formulesserrees}{%
  \setlength{\abovedisplayskip}{3pt}\setlength{\belowdisplayskip}{3pt}%
  \setlength{\abovedisplayshortskip}{2pt}\setlength{\belowdisplayshortskip}{3pt}}

\setlength{\parindent}{0pt}

\RequirePackage[implicit=false,hidelinks,pdfencoding=auto,psdextra,bookmarksopen,
  bookmarksopenlevel=1,pdfcreator={},pdfproducer={}]{hyperref}
\RequirePackage{bookmark}
\hypersetup{pdfauthor={Luc Destombe},pdfkeywords={CC BY-NC-SA 4.0}}
\newcounter{lycee@signet}
\newcommand{\signet}[3][]{\stepcounter{lycee@signet}%
  \edef\lycee@dest{\if\relax\detokenize{#1}\relax
    signet.\arabic{lycee@signet}\else#1\fi}%
  \hypertarget{\lycee@dest}{}\bookmark[level=#2,dest=\lycee@dest]{#3}}
\pdfstringdefDisableCommands{%
  \def\R{R}\def\vect#1{#1}\def\Cf{Cf}\def\Cg{Cg}\def\fois#1{×#1}%
  \def\barre#1{#1}\def\intervalle#1#2{[#1;#2]}\def\intff#1#2{[#1;#2]}%
  \def\textsuperscript#1{#1}\def\dfrac#1#2{#1/#2}\def\frac#1#2{#1/#2}%
  \def\vec#1{#1⃗}}%
\RequirePackage[official]{eurosym}

\geometry{top=0.8cm,bottom=1.3cm,left=1cm,right=1cm,footskip=0.6cm}

\pagestyle{fancy}
\fancyhf{}
\renewcommand{\headrulewidth}{0pt}
\fancyfoot[C]{\footnotesize Page \thepage{} / \pageref{LastPage}}

\newcommand{\titrefiche}[2]{%
  \begin{center}
    \Large\textbf{#1}\\[2pt]
    \large\textbf{#2}\\[2pt]
  \end{center}
  \vspace{0.10cm}}

\setlength{\columnseprule}{0.4pt}
\setlength{\columnsep}{15pt}
\renewcommand{\columnseprulecolor}{\color{black!45}}
\setlength{\multicolsep}{2pt}

\newcounter{partiefiche}
\renewcommand{\thepartiefiche}{\Roman{partiefiche}}
\newif\ifapresbandeau
\newcommand{\lycee@bandeau}[1]{%
  \par\penalty-600
  \vspace{0.08cm}%
  \noindent\colorbox{coulPart}{%
    \parbox{\dimexpr\linewidth-2\fboxsep\relax}{%
      \color{white}\bfseries\raggedright\strut#1\strut}}%
  \nopagebreak\par\nopagebreak
  \vspace{0.06cm}\nopagebreak
  \apresbandeautrue}
\newcommand{\partiefiche}[1]{%
  \refstepcounter{partiefiche}%
  \lycee@bandeau{\signet{1}{\thepartiefiche{} --- #1}\thepartiefiche{} --- #1}}
\newcommand{\partiefichesansnum}[1]{\lycee@bandeau{\signet{1}{#1}#1}}

\newcounter{exo}
\newlength{\espaceexo}\setlength{\espaceexo}{7pt}
\newlength{\espacetitre}\setlength{\espacetitre}{2pt}
\newcommand{\exo}[1][\relax]{%
  \par
  \ifapresbandeau\apresbandeaufalse\else\penalty-150\addvspace{\espaceexo}\fi
  \refstepcounter{exo}%
  \noindent\signet[exercice-\arabic{exo}]{2}{Exercice \arabic{exo}}%
  \textbf{\color{coulExo}\underline{Exercice \theexo}}%
  \ifx\relax#1\relax\else\textbf{ : #1}\fi
  \par\nopagebreak\vspace{\espacetitre}\nopagebreak
  \ignorespaces}

\setlist[enumerate]{nosep,leftmargin=*,itemsep=2pt,topsep=2pt}
\setlist[itemize]{nosep,leftmargin=*,topsep=2pt}
\newcommand{\choix}[3]{\par\hspace*{1em}%
  \makebox[0.3\linewidth][l]{a) #1}\makebox[0.3\linewidth][l]{b) #2}c) #3}
\newcommand{\deux}[2]{\par\hspace*{0.6em}%
  \makebox[0.48\linewidth][l]{#1}#2}
\newcommand{\trois}[3]{\par\hspace*{0.6em}%
  \makebox[0.32\linewidth][l]{#1}\makebox[0.32\linewidth][l]{#2}#3}

  \renewcommand{\exo}[1][\relax]{%
    \ifapresbandeau\apresbandeaufalse\else\par\penalty-150\fi
    \refstepcounter{exo}%
    \vspace{0.05cm}%
    \noindent\signet[exercice-\arabic{exo}]{2}{Exercice \arabic{exo}}%
    \textbf{\color{coulExo}\underline{Exercice \theexo}}%
    \ifx\relax#1\relax\else\textbf{ : #1}\fi%
    \par\nopagebreak
    \ignorespaces}
  \newcommand{\rappel}[1]{%
    \par\vspace{0.06cm}\nopagebreak
    \noindent\fcolorbox{black!35}{black!5}{%
      \parbox{\dimexpr\linewidth-2\fboxsep-2\fboxrule\relax}{%
        \small\textbf{Rappel.} #1}}%
    \par\vspace{0.06cm}\nopagebreak}
  \setlist[enumerate]{nosep,leftmargin=*}
  \setlist[itemize]{nosep,leftmargin=*}
  \setlength{\multicolsep}{3pt plus 1pt minus 1pt}
  \setlength{\premulticols}{10pt}
  \setlength{\postmulticols}{10pt}
  \newlist{expr}{enumerate}{1}
  \setlist[expr]{label={\Alph*\ $=$},nosep,leftmargin=*,itemsep=0pt}

\renewenvironment{center}{\par\addvspace{3pt}\centering}{\par\addvspace{3pt}}
\formulesserrees
\AtBeginDocument{\formulesserrees}
\clubpenalty=10000
\widowpenalty=10000
\displaywidowpenalty=10000
\brokenpenalty=10000
\setlength{\parskip}{0pt}

\RequirePackage{ragged2e}
\setlength{\RaggedRightRightskip}{0pt plus 1fil}
\AtBeginDocument{\RaggedRight\binoppenalty=10000 \relpenalty=10000 }
\g@addto@macro\@parboxrestore{\RaggedRight}
\makeatother

\begin{document}

\titrefiche{1\textsuperscript{re} STI2D --- Calcul littéral}{Fiche d'exercices du chapitre}

\begin{multicols}{2}\raggedcolumns

\partiefiche{Réduire une expression littérale}

\exo[Réduire une somme]
Réduire, si c'est possible, chaque expression.

\textbf{a)}
\begin{multicols}{2}
\begin{expr}
  \item $4a+3+7a$
  \item $9a+2a+5a$
  \item $6+2a+8a$
  \item $3a^{2}+7a^{2}+5a^{2}$
  \item $2a^{2}+9+4a^{2}$
  \item $3a+6a^{2}+2a^{2}$
\end{expr}
\end{multicols}

\smallskip
\textbf{b)}
\begin{multicols}{2}
\begin{expr}[start=7]
  \item $7m-12m+3m$
  \item $6m-9m+4m$
  \item $-5m-4m-6m$
  \item $9m^{2}+2m^{2}-7m^{2}$
  \item $5m-8m^{2}+3m^{2}$
  \item $-4m^{2}+6m^{2}-9m^{2}$
\end{expr}
\end{multicols}

\smallskip
\textbf{c)} Réduire, puis ranger les termes suivant les puissances décroissantes de $n$.
\begin{expr}[start=13]
  \item $3n^{2}+4n-6+5n^{2}+1+8$
  \item $-7n^{2}+9-5n+3n^{2}-n-4$
  \item $-6n^{2}-3+8-2n^{2}+5n-9n$
  \item $4+6n-3n^{2}+7n^{2}+2n-4$
\end{expr}

\exo[Réduire un produit]
Réduire chaque produit.

\textbf{a)}
\begin{multicols}{3}
\begin{expr}
  \item $8\times 2x$
  \item $3x\times 4x$
  \item $5\times 2x^{2}$
  \item $7x\times 3$
  \item $3x^{2}\times 6$
  \item $4x\times 5x$
\end{expr}
\end{multicols}

\smallskip
\textbf{b)} \textit{Attention à la règle des signes.}
\begin{multicols}{2}
\begin{expr}[start=7]
  \item $-2x\times 6x$
  \item $-9\times 3x$
  \item $5x\times(-2x)$
  \item $-3x\times(-7x)$
  \item $4\times\left(-3x^{2}\right)$
  \item $-5x^{2}\times 3$
  \item $-6\times\left(-2x^{2}\right)$
  \item $3x^{2}\times(-4)$
  \item $-7x\times 2x$
\end{expr}
\end{multicols}

\exo[Somme ou produit ?]
Réduire, si c'est possible, chaque expression. Laquelle ne se réduit pas ? Pourquoi ?
\begin{multicols}{2}
\begin{expr}
  \item $5+4x$
  \item $5\times 4x$
  \item $5x\times 4x$
  \item $5x+4x$
\end{expr}
\end{multicols}

\partiefiche{Développer}

\exo[Distributivité simple]
Développer et réduire.

\textbf{a)}
\begin{multicols}{2}
\begin{expr}
  \item $4(3x+5)$
  \item $(6+5x)\times 2$
  \item $x(2x+9)$
  \item $3x(4x+5)$
\end{expr}
\end{multicols}

\smallskip
\textbf{b)}
\begin{expr}[start=5]
  \item $5x+3(2x+4)$
  \item $2(5x+6)+4(2x+3)$
  \item $6(3+2x)+5x$
  \item $3(4x+1)+2(7+3x)$
  \item $x(3x+5)+2(6x+1)$
  \item $2x(4x+3)+3x(2x+5)$
\end{expr}

\exo[Signes et facteurs négatifs]
Développer et réduire.

\textbf{a)}
\begin{expr}
  \item $5(a+3)-4a$
  \item $-7a+3(2a-4)$
  \item $4(2+3a)-3(5+2a)$
  \item $-5(-3a+2)+2(4a-3)$
\end{expr}

\smallskip
\textbf{b)}
\begin{expr}[start=5]
  \item $3b+2b(5+6b)$
  \item $4b(-3b+5)+7b^{2}$
  \item $3b(-4b-2)-b(2b-9)$
  \item $-2b(5b-3)+6b(-2b-4)$
\end{expr}

\exo[Double distributivité]
Développer et réduire.
\begin{multicols}{2}
\begin{expr}
  \item $(2x+5)(4x+3)$
  \item $(3+4x)(5x+2)$
  \item $(5x+2)(4+3x)$
  \item $(6x+1)(3+5x)$
  \item $6x+5(2+3x)$
  \item $(4x+3)(4x+3)$
\end{expr}
\end{multicols}
\textit{Pourquoi E se développe-t-il beaucoup plus vite que les autres ?}

\partiefiche{Les trois identités remarquables}

\exo[Identité $(a+b)^{2}$]
Développer.
\begin{multicols}{2}
\begin{expr}
  \item $(3x+2)^{2}$
  \item $(4+3x)^{2}$
  \item $(2a+7)^{2}$
  \item $(5+4b)^{2}$
  \item $(x+9)^{2}$
  \item $(6x+1)^{2}$
  \item $(2m+5)^{2}$
  \item $(3+7r)^{2}$
\end{expr}
\end{multicols}

\exo[Identité $(a-b)^{2}$]
Développer.
\begin{multicols}{2}
\begin{expr}
  \item $(2x-3)^{2}$
  \item $(5x-4)^{2}$
  \item $(4-3x)^{2}$
  \item $(1-6x)^{2}$
  \item $(3a-5)^{2}$
  \item $(2-7b)^{2}$
  \item $(4m-1)^{2}$
  \item $(n-8)^{2}$
\end{expr}
\end{multicols}

\exo[Identité $(a+b)(a-b)$]
Développer. \textit{Une de ces expressions n'est pas une identité remarquable :
laquelle ?}
\begin{multicols}{2}
\begin{expr}
  \item $(3x+7)(3x-7)$
  \item $(2x+3)(2x-5)$
  \item $(5+2x)(5-2x)$
  \item $(4a+1)(4a-1)$
  \item $(b-6)(b+6)$
  \item $(9+c)(9-c)$
\end{expr}
\end{multicols}

\exo[Choisir la bonne méthode]
\textbf{a)} Développer et réduire, en choisissant à chaque fois la méthode la plus rapide.
\begin{multicols}{2}
\begin{expr}
  \item $3(2x+7)$
  \item $(2x+7)^{2}$
  \item $(5x-3)^{2}$
  \item $(7x-1)(7x+1)$
  \item $(6x+5)(2x-3)$
\end{expr}
\end{multicols}

\smallskip
\textbf{b)} Développer, réduire et ordonner.
\begin{expr}[start=6]
  \item $(3x+2)(x+4)-(2x-5)(x+1)$
  \item $(x-3)^{2}+(2x+1)(2x-1)$
  \item $(2x+5)^{2}-(x-4)(3x+2)$
\end{expr}

\partiefiche{Factoriser}

\exo[Facteur commun]
Factoriser.
\begin{multicols}{2}
\begin{expr}
  \item $6a+6b$
  \item $48a-36b$
  \item $45a-15b$
  \item $120a-80b$
  \item $7x^{2}+4x$
  \item $18x^{2}-11x$
  \item $15x+8x^{2}$
  \item $10x^{2}-4x$
  \item $4x^{2}+x$
  \item $3x^{3}-x^{2}$
\end{expr}
\end{multicols}

\exo[Identité $a^{2}+2ab+b^{2}$]
Factoriser, si c'est possible. \textit{Contrôler le double produit : deux de ces
expressions ne sont pas des carrés.}
\begin{multicols}{2}
\begin{expr}
  \item $9x^{2}+12x+4$
  \item $25x^{2}+40x+16$
  \item $9a^{2}+6a+4$
  \item $25+30b+9b^{2}$
  \item $16c^{2}+20c+25$
  \item $28d+49d^{2}+4$
\end{expr}
\end{multicols}

\exo[Identités $a^{2}-2ab+b^{2}$ et $a^{2}-b^{2}$]
\textbf{a)} Factoriser, si c'est possible. \textit{Deux de ces expressions ne sont pas
des carrés.}
\begin{multicols}{2}
\begin{expr}
  \item $16x^{2}-24x+9$
  \item $9-12x+4x^{2}$
  \item $4x^{2}-6x+9$
  \item $25x^{2}+10x-1$
  \item $9a^{2}+25-30a$
  \item $49b^{2}+4-28b$
\end{expr}
\end{multicols}

\smallskip
\textbf{b)} Factoriser à l'aide de l'identité $a^{2}-b^{2}=(a-b)(a+b)$.
\begin{multicols}{2}
\begin{expr}[start=7]
  \item $49x^{2}-25$
  \item $64-9x^{2}$
  \item $(x+2)^{2}-16$
  \item $36-(x-1)^{2}$
\end{expr}
\end{multicols}

\exo[Facteur commun entre parenthèses]
Factoriser, puis réduire chaque facteur.
\begin{expr}
  \item $(3x+4)(2x-5)+(3x+4)(5x+1)$
  \item $(6x-1)(2x+3)+(6x-1)(x-5)$
  \item $(2x+9)(4x-1)+(4x-1)(3x+2)$
  \item $(4r-3)(2r+5)+(r+7)(2r+5)$
  \item $(5x+3)(x-4)+(5x+3)$
  \item $(7r-2)(3r+1)+(7r-2)^{2}$
  \item $(4m-5)^{2}-(4m-5)(m+3)$
\end{expr}
\textit{Pour E, écrire $(5x+3)=(5x+3)\times 1$ ; pour F et G, écrire le carré comme un
produit de deux facteurs égaux.}

\partiefiche{Résoudre des équations}

\exo[Équations du premier degré]
Résoudre dans $\R$ :
\begin{multicols}{2}
\begin{enumerate}[label=\alph*)]
  \item $4x-9=3x+12$
  \item $5x-17=2x-5$
  \item $2x-11=5x+13$
  \item $3x+8=3x-4$
\end{enumerate}
\end{multicols}
\textit{Que se passe-t-il pour l'équation d) ?}

\exo[Équations avec des parenthèses]
Résoudre dans $\R$ :
\begin{enumerate}[label=\alph*)]
  \item $3(x-4)=2(-x+4)$
  \item $-2(-a+5)=4(3-a)+8$
  \item $m-(7-2m)=2m+5$
  \item $3(p-2)+4(p-5)=0$
\end{enumerate}

\exo[Équations produit]
\rappel{Un produit est nul si et seulement si l'un de ses facteurs est nul. On ne
divise \textbf{jamais} par $x$ : on factorise.}
Résoudre dans $\R$ :
\begin{multicols}{2}
\begin{enumerate}[label=\alph*)]
  \item $(3x-9)(x+5)=0$
  \item $(2x+10)(3-4x)=0$
  \item $4x(3x-8)=0$
  \item $(2x+5)^{2}=0$
  \item $6x^{2}-5x=0$
  \item $x^{2}-49=0$
\end{enumerate}
\end{multicols}
\begin{enumerate}[label=\alph*),start=7]
  \item $(x-3)(2x+5)+(x-3)(x-8)=0$
  \item $(3x+1)(x+2)-(3x+1)(2x+6)=0$
\end{enumerate}

\exo[Un premier problème]
Un magasin d'électronique reçoit $96$ composants : des résistances, des condensateurs et
des diodes. Il y a deux fois plus de résistances que de condensateurs, et $8$
condensateurs de plus que de diodes.

Combien le magasin a-t-il reçu de composants de chaque sorte ?

\partiefiche{Mise en équation}

\rappel{Inconnue (ce que représente $x$), équation, résolution, vérification, phrase de
conclusion.}

\exo[Deux commandes de câbles]
Un club de robotique achète $x$ câbles à $4$~€ l'un. Un second club achète $15$ câbles de
plus, à $1{,}50$~€ l'un. Les deux clubs dépensent la même somme.
\begin{enumerate}[label=\alph*)]
  \item Exprimer en fonction de $x$ la dépense du premier club, le nombre de câbles
        achetés par le second club, puis sa dépense.
  \item Combien de câbles le premier club a-t-il achetés ?
\end{enumerate}

\exo[Une question d'âges]
Une technicienne a $58$ ans, sa fille $33$ ans et son petit-fils $7$ ans. Elle affirme :
« Avant mes $70$ ans, la somme de vos deux âges sera égale au mien. » A-t-elle raison ?

\textit{On peut appeler $x$ le nombre d'années qui passent, mettre l'énoncé en équation,
puis comparer l'âge atteint à $70$ ans.}

\exo[Trottinettes en libre-service]
Un opérateur de trottinettes facture $1$~€ le déverrouillage, puis $0{,}25$~€ la minute.
\begin{enumerate}[label=\alph*)]
  \item Combien coûte un trajet de $12$ minutes ?
  \item Un trajet a coûté $6{,}50$~€. Combien de minutes a-t-il duré ?
  \item Un second opérateur facture $0{,}35$~€ la minute, sans déverrouillage. Pour quelle
        durée les deux opérateurs coûtent-ils le même prix ?
\end{enumerate}

\partiefiche{Résoudre des inéquations}

\rappel{Multiplier ou diviser les deux membres par un nombre \textbf{strictement
négatif} \textbf{change le sens} de l'inégalité.}

\exo[Inéquations immédiates]
Résoudre dans $\R$ et donner les solutions sous forme d'intervalle :
\begin{multicols}{2}
\begin{enumerate}[label=\alph*)]
  \item $x-5\le 0$
  \item $x+3>0$
  \item $4x\le 20$
  \item $2x+9>0$
  \item $-3x>12$
  \item $-x+7\ge 0$
  \item $5x-2<0$
  \item $\dfrac{3-4x}{5}\ge 0$
\end{enumerate}
\end{multicols}

\exo[L'inconnue des deux côtés]
Résoudre dans $\R$ :
\begin{multicols}{2}
\begin{enumerate}[label=\alph*)]
  \item $4x-5<3-2x$
  \item $3x-2\ge 6x+7$
  \item $3(x-2)\ge 9-2x$
  \item $-2(x-3)<4-x$
  \item $3(x+1)<5+3x$
  \item $2(x-4)\ge 2x+1$
\end{enumerate}
\end{multicols}
\textit{Que remarque-t-on pour e) et f) ?}

\exo[Inéquations avec des fractions]
Résoudre dans $\R$ :
\begin{multicols}{2}
\begin{enumerate}[label=\alph*)]
  \item $\dfrac{x}{2}+3\le 5$
  \item $\dfrac{3x-2}{4}<4$
  \item $\dfrac{x-1}{2}-\dfrac{2-x}{3}\ge 0$
  \item $\dfrac{x}{4}-\dfrac{6-x}{2}>3$
\end{enumerate}
\end{multicols}

\partiefiche{Tableaux de signes}

\exo[Signe d'une expression affine]
\rappel{$ax+b$ (avec $a\ne 0$) s'annule pour $x=-\dfrac{b}{a}$ ; si $a>0$, son signe est
$-$ puis $+$ ; si $a<0$, il est $+$ puis $-$.}
Dresser le tableau de signes de chaque expression.
\begin{multicols}{3}
\begin{enumerate}[label=\arabic*)]
  \item $3x-12$
  \item $-2x+10$
  \item $x+5$
  \item $7-x$
  \item $5x+2$
  \item $-4x-6$
\end{enumerate}
\end{multicols}

\exo[Signe d'un produit de deux facteurs]
Dresser le tableau de signes de chaque expression.
\begin{multicols}{2}
\begin{enumerate}[label=\arabic*)]
  \item $(x-5)(x-2)$
  \item $(1-3x)(x+4)$
  \item $(2x+6)(5-x)$
  \item $x^{2}-16$
\end{enumerate}
\end{multicols}
\textit{Pour 4), factoriser d'abord.}

\exo[Signe d'un produit de trois facteurs]
Dresser le tableau de signes de chaque expression : une ligne par facteur, puis la
ligne du produit.
\begin{enumerate}[label=\arabic*)]
  \item $x(x-2)(x+3)$
  \item $3x(2x-5)(x+4)$
  \item $\left(4-x^{2}\right)(x-3)$ \quad \textit{(factoriser d'abord $4-x^{2}$)}
\end{enumerate}

\partiefiche{Inéquations produit}

\exo[Deux facteurs]
Résoudre dans $\R$, à l'aide d'un tableau de signes :
\begin{multicols}{2}
\begin{enumerate}[label=\alph*)]
  \item $x(x-3)\ge 0$
  \item $(x+2)(x-6)<0$
  \item $(5-x)(3x+9)\ge 0$
  \item $(3x-4)(1-5x)<0$
\end{enumerate}
\end{multicols}

\exo[Trois facteurs]
Résoudre dans $\R$ :
\begin{enumerate}[label=\alph*)]
  \item $x(x-2)(x+4)\ge 0$
  \item $(x+3)(2-x)(2x-1)<0$
  \item $2x(3x-4)(x+6)\le 0$
  \item $x\left(x^{2}-16\right)>0$
\end{enumerate}

\exo[Factoriser avant de résoudre]
Résoudre dans $\R$ :
\begin{multicols}{2}
\begin{enumerate}[label=\alph*)]
  \item $x^{2}-36<0$
  \item $x^{2}-8x\le 0$
  \item $\left(9x^{2}-4\right)(x+2)>0$
  \item $(x+2)^{2}-16>0$
\end{enumerate}
\end{multicols}

\partiefiche{Inéquations quotient}

\rappel{D'abord la valeur interdite : elle annule le dénominateur, n'est jamais solution
et se marque d'une double barre dans le tableau.}

\exo[Premières inéquations quotient]
\textbf{1)} Donner la valeur interdite de
$\dfrac{4-x}{x+5}$, \ $\dfrac{7-2x}{2-x}$ \ et \ $\dfrac{x+2}{3x-9}$.

\smallskip
\textbf{2)} Résoudre dans $\R$ :
\begin{multicols}{2}
\begin{enumerate}[label=\alph*)]
  \item $\dfrac{4-x}{x+5}>0$
  \item $\dfrac{x+2}{3x-9}\le 0$
  \item $\dfrac{7-2x}{2-x}\ge 0$
  \item $\dfrac{3x+2}{4-x}<0$
\end{enumerate}
\end{multicols}

\exo[Avec une factorisation]
Résoudre dans $\R$, après avoir donné les valeurs interdites :
\begin{multicols}{3}
\begin{enumerate}[label=\alph*)]
  \item $\dfrac{x(x-2)}{5-2x}\le 0$
  \item $\dfrac{x^{2}-16}{2-x}>0$
  \item $\dfrac{6-3x}{x^{2}-9}\le 0$
\end{enumerate}
\end{multicols}

\exo[Tout ramener d'un même côté]
\rappel{On ne « fait pas passer » le dénominateur : on met tout dans le membre de
gauche, puis on réduit au même dénominateur pour comparer un seul quotient à $0$.}
Résoudre dans $\R$ :
\begin{multicols}{3}
\begin{enumerate}[label=\alph*)]
  \item $\dfrac{3x+1}{x+3}\le 2$
  \item $\dfrac{2-5x}{2-x}\ge 3$
  \item $\dfrac{x+4}{3-2x}<1$
\end{enumerate}
\end{multicols}

\partiefiche{Problèmes avec des inéquations}

\exo[Seuil de rentabilité]
Un atelier assemble des modules Bluetooth. Pour $x$ modules par jour (avec
$0<x\le 150$), le coût de production est $C(x)=15x+480$ et la recette $R(x)=31x$, en
euros.
\begin{enumerate}[label=\arabic*)]
  \item Calculer $C(25)$ et $R(25)$. L'atelier fait-il un bénéfice ce jour-là ?
  \item Résoudre $R(x)>C(x)$. À partir de combien de modules par jour l'atelier
        fait-il un bénéfice ?
  \item Le coût moyen d'un module est $\dfrac{C(x)}{x}$. Combien de modules faut-il
        assembler pour que ce coût moyen ne dépasse pas $21$~€ ?
\end{enumerate}

\exo[Louer un utilitaire électrique]
Une entreprise de livraison loue un utilitaire électrique pour la journée et parcourt
$x$ kilomètres, avec $0\le x\le 300$. Trois loueurs :
\begin{itemize}
  \item \textbf{A} : $0{,}50$~€ par kilomètre, sans forfait ;
  \item \textbf{B} : forfait de $36$~€, puis $0{,}20$~€ par kilomètre ;
  \item \textbf{C} : forfait de $60$~€, puis $0{,}10$~€ par kilomètre.
\end{itemize}
\begin{enumerate}[label=\arabic*)]
  \item Calculer le prix de $60$~km, puis de $150$~km, chez A et chez B. Quel loueur
        est le moins cher dans chaque cas ?
  \item Exprimer en fonction de $x$ le prix $A(x)$ chez A et le prix $B(x)$ chez B.
  \item Résoudre $B(x)<A(x)$ et interpréter le résultat.
  \item Chez B, l'entreprise ne veut pas dépenser plus de $80$~€. Combien de
        kilomètres peut-elle parcourir au plus ?
  \item Pour quelles distances le loueur C est-il strictement moins cher que les deux
        autres ? Justifier.
\end{enumerate}

\exo[Bénéfice d'un fabricant de trottinettes]
Un atelier fabrique chaque semaine entre $0$ et $120$ trottinettes. Pour $x$
trottinettes, son bénéfice, en euros, est
\[B(x)=-2x^{2}+240x-4\,000.\]
\begin{enumerate}[label=\arabic*)]
  \item Calculer $B(10)$ et $B(50)$. L'atelier fait-il un bénéfice dans chaque cas ?
  \item Développer $(2x-40)(100-x)$. Qu'en déduit-on pour $B(x)$ ?
  \item Résoudre $B(x)>0$ sur $\intff{0}{120}$, à l'aide d'un tableau de signes.
  \item Combien de trottinettes l'atelier doit-il fabriquer par semaine pour faire un
        bénéfice ?
  \item L'atelier vise un bénéfice d'au moins $2\,400$~€.
        Vérifier que $B(x)-2\,400=-2(x-40)(x-80)$, puis en déduire le nombre de
        trottinettes à fabriquer.
\end{enumerate}

\end{multicols}

\end{document}
